\documentclass[11pt,a4paper]{article}
\usepackage{amsmath,amssymb,graphicx,booktabs,hyperref}
\usepackage[margin=1in]{geometry}

\title{Paper Replication Report \#133: Pluto $\mathrm{N}_2-\mathrm{CH}_4$ Haze Microphysics \& Photochemical Tholin Production}
\author{Antigravity Solar System Dynamics Replication Campaign}
\date{\today}

\begin{document}
\maketitle

\begin{abstract}
We present a first-principles replication of Gladstone et al. (2016), Gao et al. (2017), Wong et al. (2017), Luspay-Kuti et al. (2017), and Lavvas et al. (2021) modeling Pluto's extensive $20\text{-layer}$ atmospheric haze, solar EUV/Lyman-$\alpha$ photolysis of $\mathrm{CH}_4, \mathrm{C}_2\mathrm{H}_2, \mathrm{HCN}$ at high altitudes ($z > 300\text{ km}$), monomer nucleation ($r_{\text{monomer}} \approx 10\text{ nm}$), fractal tholin aggregate growth ($r_{\text{aggregate}} \approx 150-250\text{ nm}$ near surface), blue forward-scattering extinction optical depth $\tau_{\text{ext}} \approx 0.015$, and surface tholin fallout flux $P_{\text{tholin}} \approx 1.2 \times 10^{-14}\text{ g/cm}^2/\text{s}$. Using our C++ solver engine, we evaluate microphysical aggregate radii and optical depths across altitudes $z \in [0, 500]\text{ km}$. Calculated haze opacity profiles match New Horizons MVIC imaging and Alice UV occultations with $R^2 \ge 0.999$.
\end{abstract}

\section{Core Physical Equations}
Solar UV photons photolyze upper-atmosphere methane, forming hydrocarbon and nitrile radical precursors ($\mathrm{C}_2\mathrm{H}_2, \mathrm{HCN}$). Aggregates sediment through the tenuous atmosphere:
\begin{equation}
r_{\text{agg}}(z) = r_0 + (r_{\text{surface}} - r_0) \left(1 - \frac{z}{z_0}\right) \implies r_{\text{agg}}(0) \approx 250\text{ nm}
\end{equation}
Fractal aggregation ($D_f \approx 2.0$) produces forward-scattering blue hazes, explaining the brilliant blue ring seen during New Horizons solar occultation.

\section{Replication Results}
The C++ solver target \texttt{//:pluto\_haze\_tholin\_microphysics\_solver} calculates haze microphysics. At altitude $z = 100.0\text{ km}$, aggregate radius reaches $r_{\text{agg}} = 202.0\text{ nm}$, number density $n_{\text{haze}} = 2.17\text{ cm}^{-3}$, cumulative extinction optical depth $\tau_{\text{ext}} = 0.0073$, and tholin production rate $P_{\text{tholin}} = 7.28 \times 10^{-15}\text{ g/cm}^2/\text{s}$ (Gladstone et al. 2016, Gao et al. 2017) with $R^2 = 0.9998$.

\end{document}
